\section{Fundamentação}

\textbf{K-means.} O K-means~\cite{lloyd:82,macqueen:67} divide $n$ exemplos em $k$ grupos, com $k$ escolhido por quem o usa. Cada grupo é representado por um centro $\mu_j$, a média dos seus exemplos, e o algoritmo busca a partição de menor \emph{inércia}, a soma das distâncias ao quadrado de cada exemplo ao centro do seu grupo:
\begin{equation*}
  \text{inércia} = \sum_{i=1}^{n} \lVert x_i - \mu_{c(i)} \rVert^2 ,
\end{equation*}
em que $c(i)$ é o grupo do exemplo $i$. O algoritmo de Lloyd começa com $k$ centros e alterna dois passos até a partição parar de mudar: cada exemplo vai para o centro mais próximo, e cada centro passa a ser a média dos exemplos do seu grupo. Nenhum dos passos aumenta a inércia, de modo que o algoritmo sempre para, mas num mínimo local, que depende dos centros iniciais. O k-means++~\cite{arthur:07} sorteia esses centros dando mais chance a exemplos longe dos já escolhidos, o que espalha os centros pelo espaço, e é comum repetir o algoritmo com vários sorteios e ficar com a partição de menor inércia. Cada iteração custa proporcionalmente ao número de exemplos, ao número de grupos e ao número de atributos, e o custo total cresce linearmente com a base. Como mede distâncias a um centro, o K-means favorece grupos arredondados e de tamanho parecido.

\textbf{Distância e escala.} Como tudo depende da distância, atributos em escalas diferentes precisam ser postos numa escala comum, ou o de maior amplitude domina. A \emph{padronização} (z-score) transforma cada valor em $(x - \bar{x})/s$, de modo que o atributo fique com média 0 e desvio padrão 1, e a escala vem de todos os valores. A \emph{normalização} (min-max) transforma cada valor em $(x - \min)/(\max - \min)$, levando o atributo para o intervalo de 0 a 1, e a escala vem só dos dois extremos: um valor muito alto comprime todos os outros perto de zero. Usamos a padronização. Atributos de texto não têm distância natural; o \emph{one-hot} cria uma coluna de 0 ou 1 para cada categoria, e duas pessoas com categorias diferentes passam a diferir em duas colunas.

\textbf{Avaliação de agrupamentos.} Sem rótulos, a qualidade de uma partição é medida por índices internos. A inércia sempre cai quando $k$ aumenta, e por isso se procura o ``cotovelo'', o $k$ a partir do qual ela cai pouco. A silhueta~\cite{rousseeuw:87} de um exemplo compara a distância média $a$ aos exemplos do próprio grupo com a distância média $b$ aos do grupo vizinho mais próximo, $(b - a)/\max(a, b)$: vai de $-1$ a 1, e a média sobre os exemplos é maior para grupos compactos e separados. O índice de Davies-Bouldin~\cite{davies:79} toma, para cada grupo, a pior razão entre a soma dos espalhamentos de dois grupos e a distância entre os seus centros, e faz a média; quanto menor, melhor. Para comparar duas partições, como as de duas seeds ou a de um agrupamento e a dos rótulos, usamos o índice de Rand ajustado (ARI)~\cite{hubert:85}, que conta os pares de exemplos postos juntos ou separados nas duas partições, corrigido pelo acaso: vale 1 para partições idênticas e cerca de 0 para concordância casual, sem depender do número dado a cada grupo. A \emph{pureza} de uma partição em relação aos rótulos é a fração de exemplos que acertariam se cada grupo respondesse o rótulo mais comum entre os seus exemplos.

\textbf{Visualização.} A análise de componentes principais (PCA) encontra as direções em que os dados mais variam e projeta os exemplos nelas. Projetar em duas componentes permite desenhar grupos formados em muitas colunas, mas mostra só a parte da variação que essas duas direções explicam.
